\documentclass[letterpaper]{article}

\usepackage[margin=1.5cm]{geometry}
\usepackage{amsmath,amsfonts}
\usepackage{url}
\usepackage{bm}

\newcommand{\Real}{\mathbb{R}}
\newcommand{\tr}{^\top}
\newcommand{\opt}{^\star}

\setlength{\parskip}{3mm plus1mm minus1mm}
\setlength{\parindent}{0mm}

\title{Problem Set 11}
\author{Foundations of Machine Learning}
\date{}

\begin{document}
\maketitle

\begin{itemize}
\item \textbf{Submission}: Turn in as a \texttt{PDF} and a \textbf{separate source code file} on \url{http://mycourses.unh.edu}. Scanned hand-written notes are encouraged.
\item \textbf{Questions}: See the class website for a link to office hours
\end{itemize}

% Covers: Probability basics: random variables (MML 6.1--6.3) and distributions / Gaussian (MML 6.4--6.5).

\textbf{Relevant reading}: MML 6.1--6.5. Problems marked \textbf{[code]} should be submitted as a program in a language of your choice (Julia or Python recommended). The cited exercise numbers are from MML (Deisenroth et al., 2021).

\begin{enumerate}
\item MML, Chapter 6 exercises: any three on random variables, expectation, and Bayes' theorem (e.g.\ Exercises 6.1--6.3).
\item MML, Chapter 6 exercises: any two on joint/marginal distributions and the Gaussian (e.g.\ Exercises 6.4--6.6).
\item \textbf{[code]} Draw samples from a two-dimensional Gaussian with a chosen mean and covariance. Estimate the sample mean and covariance from the draws, and produce a scatter plot with the estimated mean marked.
\item \textbf{Time spent.} Report approximately how many hours you spent completing this problem set (including the reading).
\end{enumerate}

\end{document}
